
In games with polynomial utility functions, the best response oracles can employ both methods of the methods of global optimization described in Chapter \ref{ch:global}. In the rest of the games we use the spatial Branch and Bound approach as the best response oracle. From our experience, local solvers can work sufficiently well as oracles in some games, although no convergence guarantees then apply. \todo[inline]{this belongs elsewhere (to the "this is a metaalgorithm parametrized by master and oracle...".... ok now it's here)}

blah...The examples in our collections have utilities specified in terms of sums and products of functions such as logarithms and square roots. All such functions should be accepted either directly by global solvers or through external modeling libraries such as \textsc{ampl}.
\todo[inline]{merge this properly}

We show a comparison of Oracle implementations based on non-linear solvers.
While we have succesfully used local-solvers as oracles, it is important to remember that convergence to equilibria is no longer guaranteed.
Regardless, we compare their performance in some of the following experiments.

The experiments are split into three groups, random polynomials, non-polynomial test functions, and a practical experiment
where we solve a where we run a fixed number of iterations of Double Oracle on the Colonel Blotto game.


\paragraph{polynomials}
blah

\paragraph{non-polynomial test functions}
blah

The next experiment is a measurement of the runtime of different solvers in practice over a course of ten iterations of double oracle in a game of General Blotto. The utility functions in \emph{General Blotto} games are sums of non-linearly weighted margins of victory over sets of fronts.
We chose a game with three isolated fronts defined by the utility $u_1(x_1, x_2) = \sum_{j=1}^m f\left( x_{1j} - x_{2j} \right)$ with a non-differentiable weighing function $f(z) = \operatorname{sgn}(z)\,|z|^{0.5}$. This game exhibits very slow convergence with the Double Oracle algorithm.
\iffalse
Multilinear terms, square roots and sums can be formulated directly via the generic modelling languages \emph{JuMP}, and \emph{AMPL}; however the sign function and absolute value are more challenging. Due to a lack of support and differentiability, all of the solvers: \emph{Ipopt}, \emph{Bonmin}, \emph{Couenne}, \emph{SCIP}, and \emph{Gurobi} were unable to solve a direct formulation of the oracle problem. Instead, we reformulated the weighing function $f$ using a binary variable as $f(x) = (2b-1)z$, subject to the constraints:
\[
\begin{aligned}
b &\in \{0,1\} \\
s,n,p &\in [0,1] \\
p &\leq b \\
n &\leq 1-b \\
x &= p - n \\
z^2 &=  p + n. \\
\end{aligned}
\]
The reformulation is a mixed-integer quadratic program.
\fi